\begin{abstract}
Large-scale unsupervised language models (LMs) acquire extensive world knowledge and demonstrate some capacity for reasoning; however, exerting fine-grained control over their outputs remains challenging due to their reliance on purely unsupervised training processes. To address steerability, existing approaches collect human-annotated preference data regarding model outputs and adjust the LMs to match these preferences, often leveraging reinforcement learning from human feedback (RLHF). Despite its success, RLHF involves a multi-step and frequently unstable protocol: first, constructing a reward model based on human preference annotations, then employing reinforcement learning to maximize the learned reward function during fine-tuning, all while avoiding significant deviation from the base LM. In this work, we propose a novel reward model formulation within RLHF that permits a closed-form solution for the optimal policy, reducing the standard RLHF problem to optimization with a straightforward classification objective. This approach, termed \textit{Direct Preference Optimization} (DPO), is efficient, robust, and requires considerably fewer computational resources, as it eliminates the need for language model sampling during fine-tuning and minimizes hyperparameter sensitivity. Empirical results indicate that DPO aligns language models with human preferences at least as well as, and often surpasses, prior techniques. Remarkably, DPO outperforms PPO-based RLHF methods in modulating generation sentiment and achieves comparable or superior outcomes in tasks like summarization and single-turn conversation, all while offering a much more straightforward implementation and training process.
\end{abstract}

\section{Introduction}
Large-scale unsupervised language models (LMs) trained on extensive corpora have demonstrated unexpected and remarkable abilities~\citep{chowdhery2022palm, brown2020language, touvron2023llama, bubeck2023sparks}. Nevertheless, the data used to train these models originates from human sources characterized by diverse intentions, expertise, and preferences. Many of these human behaviors may be unsuitable for direct emulation. For instance, we may want an AI code assistant to \textit{recognize} common errors in order to provide corrections, yet when generating new code, it should be biased toward the (potentially infrequent) instances of expert-level programming present in its training set. Likewise, it is beneficial for a language model to \textit{be informed of} prevalent misconceptions, such as one believed by half of users; however, it is undesirable for the model to replicate or endorse the misconception with corresponding frequency in its output. Put differently, isolating and promoting the model's \emph{appropriate behaviors and outputs} from within its expansive \textit{knowledge and repertoire} is critical for developing AI systems that are controllable, high-performing, and aligned with safety objectives~\citep{ouyang2022training}. While prevailing techniques tend to align LMs with human values through reinforcement learning (RL) based preference modeling, we will demonstrate that the objective in these RL-based preference optimization methods can be equivalently formulated and exactly optimized as a simple binary cross-entropy loss, streamlining the process of learning preferences.

Broadly speaking, contemporary approaches introduce desired model behaviors by leveraging carefully selected sets of human preference data that exemplify safety and helpfulness. This stage of learning human preferences typically follows an initial unsupervised pre-training phase utilizing a large-scale text corpus. Although one could employ supervised fine-tuning directly on human-authored exemplars of high-quality output, the dominant strategies instead involve reinforcement learning from human or artificial feedback (RLHF/RLAIF; \citep{christiano2017deep, bai2022constitutional}). In RLHF, a reward model is constructed from annotated human preferences and then RL is used to update the language model policy to favor responses that earn higher rewards while ensuring that its behavior does not deviate drastically from that of the original model. While RLHF techniques lead to LMs with notable conversational and coding skills, the RLHF workflow is substantially more intricate than conventional supervised training. It requires the training of multiple language models and in-loop policy sampling during training, resulting in considerable computational overhead.

This work introduces a method for optimizing language models to follow human preferences without relying on explicit reward modeling or reinforcement learning techniques. We present \textit{{Direct Preference Optimization} (DPO)}, a new algorithm that implicitly targets the same optimization goal as popular RLHF methods—maximizing reward while maintaining a KL-divergence constraint—but does so via a design that is both easy to implement and train. Conceptually, the DPO approach enhances the log odds of preferred outputs relative to less preferred ones, incorporating a dynamic importance weight for each data point. This mechanism guards against the model collapse commonly seen with more naive probability ratio-based objectives. In a manner similar to prior approaches, DPO assumes a theoretical framework for preferences (e.g., the Bradley-Terry model; \cite{bradley1952rankanalysis}) that evaluates the alignment of a reward function with observed human preferences. Yet, rather than employing this preference model to train a reward function and subsequently learning a policy to optimize the reward, DPO re-parameterizes the preference loss to make it a direct function of the policy itself. Using datasets of pairwise human preferences over model outputs, DPO enables optimization of a policy using a standard binary cross-entropy loss, effectively learning the policy optimal with respect to an implicit reward function consistent with the observed preferences.

The primary contribution of our work is the development of Direct Preference Optimization (DPO), an RL-free approach to preference-based language model training. Empirical results demonstrate that DPO matches or exceeds the performance of established techniques, such as PPO-based RLHF, for applications including sentiment adjustment, summarization, and conversational modeling, with models of up to 6B parameters.

Progressively larger self-supervised language models are capable of addressing certain tasks in a zero-shot setting \citep{radford2019language} or with the aid of few-shot prompts \citep{gpt3,megatron,chowdhery2022palm}. Nevertheless, their outcomes on downstream applications and their adherence to user instructions are markedly enhanced when these models are fine-tuned on corpora containing explicit instructions alongside human-authored responses \citep{mishra-etal-2022-cross,sanh2022multitask,chung2022scaling,thoppilan2022lamda}. Through this process, termed `instruction-tuning,' large language models (LLMs) demonstrate improved capacity to generalize to prompts not present in the original fine-tuning data, thus elevating their general utility \citep{chung2022scaling}. While instruction tuning has proven beneficial, the collection of human \textit{relative} assessments of output quality is often simpler than amassing expert-generated examples. Subsequent research has therefore explored the fine-tuning of LLMs using preference datasets, where models are exposed to comparisons reflecting human choices; this strategy has shown gains in translation \citep{kreutzer-etal-2018-reliability}, summarization \citep{stiennon2022learning,ziegler2020finetuning}, narrative generation \citep{ziegler2020finetuning}, and following user directions \citep{ouyang2022training,ramamurthy2023is}. Typically, such approaches begin by training a neural reward model to predict preferences—commonly employing frameworks like the Bradley-Terry model \citep{bradley1952rankanalysis}—and then adjust the language model to maximize this learned reward signal via reinforcement learning techniques, such as REINFORCE \citep{williams1992reinforce}, proximal policy optimization (PPO; \cite{schulman2017proximal}), or related variants \citep{ramamurthy2023is}. Related studies have further utilized instruction-following LLMs refined with human feedback to create additional artificial preference data targeting qualities like safety or non-maliciousness \citep{bai2022constitutional}, relying on minimally detailed human-written rubrics to steer annotation. This synthesis merges traditions in reinforcement learning-based language model training for diverse tasks~\citep{Ranzato2015SequenceLT,paulus2018a,wu2018learning} with more general approaches to preference-based learning from humans \citep{christiano2017deep,kupcsik2018learning}. Although exploiting relative human judgments is an attractive strategy, using reinforcement learning to adapt large language models is still fraught with practical hurdles; in response, this work introduces a principled methodology for preference optimization that avoids the need for reinforcement learning.

Beyond the realm of language, the problem of learning policies from preferences has garnered attention in both bandit and reinforcement learning frameworks, leading to a variety of proposed solutions. Contextual bandit algorithms that learn from ranked choices or preferences between actions, as opposed to explicit reward feedback, are referred to as contextual dueling bandits (CDB; \cite{yue2012karmed,dudik2015contextual}). In such settings, where absolute reward information is unavailable, theoretical characterizations of optimality in CDBs typically rely on the concept of a \textit{von Neumann winner}: a policy whose probability of outperforming any other given policy is no less than 50\% in expectation \citep{dudik2015contextual}. Unlike CDBs, which receive preference feedback in an online fashion, human preference learning usually proceeds from a static set of offline, preference-labeled action pairs \citep{yan2022human}. Similarly, in \textit{preference-based RL} (PbRL), agents adapt according to binary preference signals derived from an unknown 'scoring' function, instead of direct reward observations \citep{BusaFekete2014,ruiz2023dueling}. Multiple strategies have been proposed for PbRL—including those capable of leveraging off-policy data—though these typically consist of a two-step process: first, an explicit estimate of the hidden scoring or reward function is learned, and second, the policy is optimized using this estimate \citep{jain2013learning,BusaFekete2014,christiano2017deep,sadigh2017active,kupcsik2018learning}. By contrast, our approach focuses on a unified, single-phase policy learning method that directly aligns the policy with observed preferences.

\section{Preliminaries}\label{section:prelims}

We begin with an overview of the RLHF procedure outlined by \citeauthor{ziegler2020finetuning} and extended in later works \citep{stiennon2022learning, bai2022training, ouyang2022training}. This process generally unfolds in three main stages: 1) supervised fine-tuning (SFT), 2) gathering preference data and learning a reward model, and 3) reinforcement learning for policy optimization.

\textbf{SFT}: The typical RLHF pipeline initiates with the supervised fine-tuning of a pre-trained language model on curated datasets suited to the downstream application (such as dialogue or summarization tasks), resulting in an adapted model $\pi^\text{SFT}$.

\textbf{Reward Modelling Phase}: Subsequently, this fine-tuned model is prompted with input $x$ and used to generate pairs of responses $(y_1, y_2)\sim \pi^\text{SFT}(y \mid x)$. These response pairs are then evaluated by human annotators, who select their preferred answer, expressed as $y_w\succ y_l \mid x$, where $y_w$ indicates the favored completion and $y_l$ the less favored one from the pair $(y_1, y_2)$. It is assumed that such preferences are the outcome of an underlying, unobserved reward function $r^*(y, x)$. To represent human preferences, several statistical models may be employed; the Bradley-Terry (BT) model \cite{bradley1952rankanalysis} is commonly used for pairwise comparisons, although extensions to broader Plackett-Luce ranking models \citep{plackett1975analysis, luce2012individual} are feasible given multiple ranked options. Under the BT framework, the human-induced preference probability $p^*$ is characterized as follows:

Given a static set of preference data $\mathcal{D}=\bigl\{x^{(i)}, y_w^{(i)}, y_l^{(i)}\bigr\}_{i=1}^N$ collected from $p^*$, we introduce a parametric reward function $r_{\phi}(x, y)$ and estimate its parameters through the maximum likelihood approach. When framed as a binary classification task, the corresponding negative log-likelihood loss takes the following form:

\begin{equation}\label{eq:reward_model}
    \mathcal{L}_R(r_{\phi}, \mathcal{D}) = -\mathbb{E}_{(x, y_w, y_l)\sim \mathcal{D}}\bigl[\log \sigma(r_{\phi}(x, y_w)- r_{\phi}(x, y_l))\bigr]
\end{equation}

where $\sigma$ denotes the logistic function. For language model applications, $r_{\phi}(x, y)$ is commonly initialized from a supervised fine-tuning model $\pi^\text{SFT}(y \mid x)$ and supplemented with a linear layer atop the transformer's final representation to yield a scalar reward output~\cite{ziegler2020finetuning}. Previous studies have found it helpful to normalize the rewards, enforcing $\mathbb{E}_{x,y\sim \mathcal{D}}\left[r_\phi(x, y)\right] = 0$ for every $x$, in order to stabilize the learned reward values.

\textbf{RL Fine-Tuning Phase}: In the reinforcement learning phase, this reward function supplies the training signals to the language model. Consistent with existing literature~\citep{jaques2017sequence, jaques2020human}, the objective is given by:

\begin{equation}\label{eq:RL}
\max_{\pi_{\theta}}  \mathbb{E}_{x\sim \mathcal{D}, y\sim \pi_{\theta}(y \mid x)}\bigl[r_{\phi}(x, y)\bigr] - \beta\mathbb{D}_{\textrm{KL}}\bigl[\pi_{\theta}(y\mid x)\mid \mid \pi_\text{ref}(y\mid x)\bigr],
\end{equation}

where $\beta$ modulates the divergence penalty relative to the reference policy $\pi_\text{ref}$, typically selected as the initial supervised fine-tuned model $\pi^\text{SFT}$. The policy $\pi_\theta$ is generally initialized with the same weights as $\pi^\text{SFT}$. This KL constraint is crucial; it not only preserves generation diversity and mitigates the risk of mode collapse to a narrow set of high-reward outputs, but also keeps the policy distribution within the region where the reward function remains accurate. Given that text generation is inherently discrete, differentiating through this objective is infeasible, necessitating reinforcement learning techniques for optimization. The prevalent method~\citep{ziegler2020finetuning, stiennon2022learning, bai2022training, ouyang2022training} constructs the composite reward function as ${r(x, y) = r_{\phi}(x, y) -\beta (\log \pi_{\theta}(y\mid x) - \log \pi_\text{ref}(y\mid x))}$, which is then optimized by proximal policy optimization (PPO)~\cite{schulman2017proximal}. 

\section{Direct Preference Optimization}\label{sec:DPO}

Given the computational demands and complexities that come with standard RL algorithms on large-scale language model training, our aim is to develop a streamlined policy optimization method that works directly from user preference data. Rather than the two-stage approach in RLHF (learning a reward model and subsequently optimizing it with RL), we instead exploit a reward model parameterization whose optimal policy admits a closed-form solution, thereby sidestepping the need for an RL-based optimization loop. Central to our methodology is the analytical relationship mapping reward models to their corresponding optimal policies, enabling a reformulation of loss functions over rewards as loss functions over policies themselves. This variable substitution eliminates the necessity of explicitly learning a standalone reward model, yet still faithfully encodes established human preference frameworks such as the Bradley-Terry model. In this structure, the policy model jointly serves as both the language generator and an implicit reward estimator.

\textbf{Derivation of the DPO objective.} As a starting point, we consider the standard RL objective outlined in Eq.~\ref{eq:RL}, defined for an arbitrary reward function $r$. Building on prior developments in the literature~\citep{peters2007reinforcement, peng2019advantage, korbak2022reinforcement, go2023aligning}, one can directly verify that the solution which maximizes the expected reward subject to a KL-divergence constraint—see Eq.~\ref{eq:RL}—has the following parametric form:

\begin{equation}\label{eq:op_policy}
    \pi_r(y\mid x) = \frac{1}{Z(x)}\pi_\text{ref}(y\mid x)\exp\left(\frac{1}{\beta}r(x, y)\right),
\end{equation}

with the partition function given by $Z(x) =\sum_{y}\pi_\text{ref}(y\mid x)\exp\left(\frac{1}{\beta}r(x, y)\right)$. Detailed steps of this derivation appear in Appendix \ref{app:derivation1}. Despite employing $r_{\phi}$—the maximum likelihood estimator for the true reward $r^*$—computing the partition function $Z(x)$ remains computationally intensive \citep{korbak2022reinforcement, go2023aligning}, rendering the practical use of this formula challenging. Nonetheless, Eq.~\ref{eq:op_policy} can be rearranged to reformulate the reward $r$ as a function of the corresponding optimal policy $\pi_r$, reference policy $\pi_\text{ref}$, and the, in general, unknown $Z(\cdot)$. Specifically, applying the logarithm to both sides of Eq.~\ref{eq:op_policy}, and simplifying, yields:

\begin{equation}\label{eq:main_eq}
    r(x,y) =\beta \log \frac{\pi_r(y\mid x)}{\pi_\text{ref}(y\mid x)} + \beta \log Z(x).
\end{equation}

This reformulation can be used for both the actual reward function $r^*$ and its associated optimal policy $\pi^*$. Crucially, since the Bradley-Terry model only considers differences in reward between pairs of outputs—namely, ${p^*(y_1 \succ y_2 \mid x) = \sigma(r^*(x, y_1) - r^*(x, y_2))}$—one can substitute the expression for $r^*$ in Eq.~\ref{eq:main_eq} into the preference probability model (Eq.~\ref{eq:bradley-terry}). Due to this subtraction, the partition function terms $Z(x)$ cancel, which allows us to recast the preference probability using only $\pi^*$ and $\pi_\text{ref}$. Therefore, the optimal policy $\pi^*$ learned by RLHF within the Bradley-Terry framework is governed by the preference relation:

\begin{equation}\label{eq:objective}
    p^*(y_1\succ y_2 \mid x)=\frac{1}{1 + \exp\left(\beta \log \frac{\pi^*(y_2\mid x)}{\pi_\text{ref}(y_2\mid x)} - \beta \log \frac{\pi^*(y_1\mid x)}{\pi_\text{ref}(y_1\mid x)}\right)}
\end{equation}

Full details of this derivation are provided in Appendix~\ref{app:derivation2}. Although Eq.~\ref{eq:objective} is derived with the Bradley-Terry model, the same approach extends naturally to the more general Plackett-Luce models~\citep{plackett1975analysis, luce2012individual}, as detailed in Appendix~\ref{app:plackett_luce_models}.

By expressing the likelihood of observed human preference data in terms of the optimal policy (rather than the reward function), we may now pose a maximum likelihood estimation problem for a parametrized policy $\pi_\theta$. In analogy with the reward modeling objective (cf. Eq.~\ref{eq:reward_model}), this leads us to the following policy-level learning criterion:

\begin{equation}\label{eq:optimum_model}
    \mathcal{L}_\text{DPO}(\pi_{\theta}; \pi_\text{ref}) = -\mathbb{E}_{(x, y_w, y_l)\sim \mathcal{D}}\left[\log \sigma \left(\beta \log \frac{\pi

In this formulation, we learn an implicit reward through an alternative form of parameterization, such that the resulting optimal policy corresponds directly to $\pi_\theta$. Furthermore, by noting the equivalence of our method to learning a reparameterized Bradley-Terry model, we inherit specific theoretical guarantees, including consistency given appropriate assumptions about the preference data distribution \cite{bong2022generalized}. We elaborate on the theoretical implications of DPO, and its connections to related literature, in Section~\ref{sec:theory}.

\textbf{Interpretation of the DPO update.} To gain a mechanistic perspective on DPO, it is instructive to consider the gradient of its loss, $\mathcal{L}_\text{DPO}$. The derivative with respect to the parameters $\theta$ is expressed as follows:

\begin{multline*}\label{eq:gradient}
    \nabla_\theta \mathcal{L}_\text{DPO}(\pi_\theta;\pi_\text{ref}) = \\ -\beta\mathbb{E}_{(x, y_w, y_l) \sim \mathcal{D}} \bigg[\underbrace{\sigma(\hat{r}_\theta(x, y_l) - \hat{r}_\theta (x, y_w))}_\text{assigns more weight to errors in reward ranking}\bigg[\underbrace{\nabla_\theta\log \pi(y_w \mid x)}_\text{boosts likelihood of $y_w$} - \underbrace{\nabla_\theta\log\pi(y_l \mid x)}_\text{reduces likelihood of $y_l$}\bigg]\bigg],
\end{multline*}

Here, $\hat{r}_\theta(x, y) = \beta \log \frac{\pi_\theta(y \mid x)}{\pi_\text{ref}(y \mid x)}$ defines the implicit reward function, constructed from the language model $\pi_\theta$ and a fixed reference $\pi_\text{ref}$ (see additional discussion in Section~\ref{sec:theory}). Conceptually, this gradient encourages the model to assign greater probability to preferred outputs $y_w$ and suppress the probability of less favored outputs $y_l$. Crucially, each sample's contribution is modulated by how much the implicit reward model $\hat{r}_\theta$ overestimates the value of the less-preferred option $y_l$ compared to $y_w$, amplified by the factor $\beta$, thereby quantifying misorderings of the completions under the implicit reward—this also incorporates the effect of the KL regularization. Our empirical results underscore the significance of this weighting mechanism: removing it and using a uniform weighting leads to degenerative behaviors in the language model (see Appendix Table~\ref{tab:unlikelihood_generations}).

\textbf{DPO Overview.} The standard DPO workflow consists of the following steps: 1) For each prompt $x$, generate completions $y_1, y_2 \sim \pi_\text{ref}(\cdot \mid x)$ and assign human preference labels to create the offline preference dataset $\mathcal{D} = \{x^{(i)}, y_w^{(i)}, y_l^{(i)}\}_{i=1}^N$; 2) Train the language model $\pi_\theta$ by minimizing $\mathcal{L}_\text{DPO}$ given $\pi_\text{ref}$, $\mathcal{D}$, and a specified $\beta$. In real-world scenarios, preference datasets that are already available are often preferred over the costly process of collecting new samples and human annotations. Since these datasets are typically derived from $\pi^\text{SFT}$, we set $\pi_\text{ref} = \pi^\text{SFT}$ when it is accessible. When $\pi^\text{SFT}$ is missing, $\pi_\text{ref}$ is instead obtained by maximizing the likelihood of preferred completions ${(x, y_w)}$, i.e., ${\pi_\text{ref} = \argmax_{\pi}\mathbb{E}_{x, y_w \sim \mathcal{D}}\left[\log \pi(y_w \mid x)\right]}$. This strategy alleviates issues arising from a mismatch between the unknown true reference distribution and the $\pi_\text{ref}$ employed by DPO. Further implementation specifics and choices of hyperparameters are detailed in Appendix~\ref{app:implementation}.

\section{Theoretical Analysis of DPO} In what follows, we deepen our interpretation of the DPO approach, establish its theoretical foundation, and highlight how DPO addresses limitations inherent in actor-critic RLHF algorithms (including PPO~\cite{schulman2017proximal}).

\label{sec:theory}

\subsection{Your Language Model Is Secretly a Reward Model} DPO sidesteps both explicit reward model training and reinforcement learning for policy optimization by consolidating these steps into a single maximum likelihood learning procedure. Importantly, the objective in Eq. \ref{eq:main_eq} mirrors the Bradley-Terry model where rewards are parameterized as $r^*(x, y) = \beta \log\frac{\pi^*_\theta(y \mid x)}{\pi_\text{ref}(y \mid x)}$, and training the model $\pi_{\theta}$ corresponds to optimizing the reward model as in Eq. \ref{eq:reward_model}, under a change of variables. This section develops the theoretical justification for such a reparameterization, demonstrates that it does not limit the class of attainable reward models, and confirms that the optimal policy can indeed be recovered precisely. We begin by introducing an equivalence relation over reward functions.

\begin{definition}
Two reward functions $r(x, y)$ and $r'(x, y)$ are said to be equivalent if and only if ${r(x, y)-r'(x, y) = f(x)}$ for some function $f$.
\end{definition}

It is straightforward to verify that this definition forms an equivalence relation, resulting in a partitioning of the space of reward functions into equivalence classes. We can establish the following lemmas:

\begin{lemma}\label{lemma:same_prefrence} 
Within the Plackett-Luce preference model, and the Bradley-Terry model as a special case, any pair of reward functions from the same equivalence class gives rise to identical preference distributions.
\end{lemma}

\begin{lemma}\label{lemma:same_policy}
    If two reward functions belong to the same equivalence class, they will yield the same optimal policy under the constrained reinforcement learning (RL) setup.
\end{lemma}

Since their proofs are elementary, we provide details in Appendix \ref{app:lemma1}. The first lemma highlights a familiar identification problem present in the Plackett-Luce models \cite{plackett1975analysis}, arising from their inherent under-specification. Consequently, it is common practice to impose extra identifiability constraints in order to ensure meaningful maximum likelihood estimates, as discussed in Eq. \ref{eq:reward_model} \cite{bong2022generalized}. According to the second lemma, all reward functions within the same equivalence class are associated with the same optimal policy. Therefore, in our analysis, it suffices to recover any reward function from the correct class for policy optimization. We formalize this in the following theorem, whose proof is provided in Appendix~\ref{app:thm1}:

\begin{theorem}\label{thm:main}
    Given mild regularity conditions, any reward equivalence class compatible with the Plackett-Luce (including Bradley-Terry) model can be represented in the form ${r(x, y) = \beta \log \frac{\pi(y\mid x)}{\pi_\text{ref}(y\mid x)}}$ for some conditional model $\pi(y\mid x)$ and a fixed reference model $\pi_\text{ref}(y \mid x)$.
\end{theorem}

\begin{sproof}
    Take an arbitrary reward function $r(x, y)$, and let $\pi_r(y \mid x)$ be the corresponding optimal policy as defined in Eq. \ref{eq:op_policy}. We aim to show that an equivalent reward function to $r(x, y)$ can be expressed through the proposed reparameterization. Define the projection operator $f$ as
\begin{equation}
    f(r; \pi_\text{ref}, \beta)(x, y) = r(x, y) - \beta\log\sum_{y}\pi_\text{ref}(y\mid x)\exp\left(\frac{1}{\beta}r(x, y)\right)
\end{equation}
This operator normalizes the original reward by subtracting the log partition function (dependent only on $x$), so $f(r; \pi_\text{ref}, \beta)(x, y)$ remains in the same equivalence class as $r(x, y)$. Substituting the right-hand side of Eq.~\ref{eq:main_eq} (which applies to any reward function), we obtain $f(r; \pi_\text{ref}, \beta)(x, y) = \beta \log \frac{\pi_r(y\mid x)}{\pi_\text{ref}(y\mid x)}$. Thus, this projection produces a representative of the equivalence class of $r$ in the required form, implying no loss of generality by employing the suggested reparameterization.
\end{sproof}

Theorem~\ref{thm:main} can be interpreted as identifying the specific reward function that the DPO reparameterization picks within each equivalence class—namely, the function $r(x, y)$ that satisfies the following criterion:

\begin{equation}\label{eq:lag_p}
     \sum_{y}\underbrace{\pi_\text{ref}(y\mid x)\exp\left(\frac{1}{\beta}r(x, y)\right)}_{=\pi(y\mid x)\text{, using Thm.~\ref{thm:main} reparam.}} = 1,
\end{equation}

In other words, $\pi(y\mid x)$ constitutes a legitimate probability distribution (i.e., non-negative probabilities summing to unity). From Eq.~\ref{eq:op_policy}, we recognize that Eq.~\ref{eq:lag_p} represents the partition function corresponding to the optimal policy defined by the reward $r(x, y)$. The central insight offered by the DPO algorithm is that by introducing particular constraints on the generally under-determined Plackett-Luce family of preference models (including, as a special case, the Bradley-Terry model), it is possible to retain the set of expressible reward functions, while simultaneously rendering the optimal policy in Eq.~\ref{eq:op_policy} analytically manageable across all prompts $x$.

\subsection{Instability of Actor-Critic Algorithms} Our framework also sheds light on the instability issues often encountered with actor-critic approaches such as PPO, which are prevalent in RLHF setups. We limit our analysis to the RL fine-tuning stage described in Section~\ref{section:prelims}, drawing parallels to the control-as-inference formalism \cite{levine2018reinforcement} for the constrained RL formulation in \ref{eq:RL}. With a parametrized model $\pi_{\theta}(y\mid x)$, we aim to minimize the Kullback-Leibler divergence $\mathbb{D}_{\text{KL}}[\pi_{\theta}(y|x) \mid \mid \pi^*(y\mid x)]$, where the reference optimal policy $\pi^*$ arises from Eq.~\ref{eq:optimum_model}, underpinned by the reward $r_{\phi}(y, x)$. This leads, through straightforward algebra, to the following objective:

\begin{equation}\label{eq:AC}
    \max_{\pi_{\theta}}\mathbb{E}_{\pi_{\theta}(y\mid x)}\bigg[\underbrace{r_{\phi}(x, y) -\beta\log\sum_{y}\pi_\text{ref}(y\mid x)\exp\left(\frac{1}{\beta}r_{\phi}(x, y)\right)}_{f(r_{\phi}, \pi_\text{ref}, \beta)} - \underbrace{\beta\log\frac{\pi_{\theta}(y\mid x)}{\pi_\text{ref}(y\mid x)}}_{\text{KL}}\bigg]
\end{equation}

This objective aligns with those pursued in previous studies 
\citep{ziegler2020finetuning, stiennon2022learning, bai2022training, ouyang2022training}, where the DPO-equivalent reward is employed for the reward class $r_{\phi}$. Here, the normalization component in $f(r_{\phi}, \pi_\text{ref}, \beta)$ can be viewed as representing the soft value function corresponding to the reference policy $\pi_\text{ref}$. Although this normalization does not influence the set of optimal policies, its absence could result in the objective’s policy gradient exhibiting high variance, thereby destabilizing training. One way to handle this normalization is by introducing a learned value function, though optimizing such a function can itself pose challenges. Another approach, taken in earlier work, involves normalizing the reward using a baseline derived from human completions—essentially approximating the normalization with a single-sample Monte Carlo estimate. Distinctly, the DPO reparameterization offers a reward formulation that obviates the need for any explicit baseline. 

\section{Experiments}
This section presents empirical analyses of DPO’s capacity to learn policies based solely on preference data. We begin by examining a controlled text-generation scenario, probing how effectively DPO balances reward maximization against maintaining low KL-divergence from the reference policy, in comparison to widely used preference learning techniques such as PPO. Subsequently, we assess DPO on larger-scale models and challenging RLHF benchmarks, including tasks in summarization and dialogue. Notably, our results suggest that with minimal hyperparameter tuning, DPO matches or exceeds the performance of robust baselines, including RLHF with PPO and selecting the best out of $N$ trajectories according to a learned reward. Preceding the presentation of these empirical findings, we outline our experimental methodology, with further specifics provided in Appendix~\ref{app:exp_details}.

\textbf{Tasks.} We conduct our study across three distinct open-domain text generation settings. In each case, the learning algorithms derive a policy from a preference dataset, denoted as $\mathcal{D}=\bigl\{x^{(i)}, y_w^{(i)}, y_l^{(i)}\bigr\}_{i=1}^N$. In the \textbf{controlled sentiment generation} task, $x$ corresponds to an initial fragment of a movie review sourced from the IMDb dataset \cite{maas-EtAl:2011:ACL-HLT2011}; the objective for the policy is to extend $x$ by generating $y$ that exhibits positive sentiment. To achieve a controlled assessment, we construct preference pairs automatically by leveraging a pre-trained sentiment classifier to ensure $p(\text{positive}\mid x,y_w)>p(\text{positive}\mid x,y_l)$. For supervised fine-tuning (SFT), we adapt GPT-2-large using the training split of IMDB reviews until satisfactory convergence (for additional methodology, see App~\ref{app:sentiment_details}). For the \textbf{summarization} task, $x$ is taken as a Reddit forum post, and the policy's role is to produce a summary $y$ that distills the primary points. We adopt the Reddit TL;DR summarization dataset \citep{volske-etal-2017-tl} and leverage human preference annotations from \citeauthor{stiennon2022learning}, consistent with existing literature. Our SFT baseline is a model further trained on summaries authored by humans,\footnote{\url{https://huggingface.co/CarperAI/openai_summarize_tldr_sft}} utilizing the TRLX framework \citep{leandro_von_werra_2023_7790115} for RLHF. The human preferences were originally elicited on completions from an alternative but comparably fine-tuned SFT model by \citeauthor{stiennon2022learning}. In the \textbf{single-turn dialogue} setting, $x$ is a single user input, which can range from complex scientific questions to personal advice requests. Here, the policy is required to generate a relevant and informative reply $y$. We employ the Anthropic Helpful and Harmless dialogue dataset \citep{bai2022training}, comprising 170,000 dialogue exchanges between humans and an automated system, where each conversation concludes with two candidate model responses and an annotation identifying which one is preferred by humans. In the absence of a publicly available pre-trained SFT model for this task, we create one by fine-tuning a general-purpose language model solely on human-preferred responses.

\textbf{Evaluation.} We adopt two distinct evaluation strategies in our experiments. To assess the performance of each algorithm with respect to the constrained reward maximization goal, we examine their respective trade-off frontiers between attained reward and KL-divergence from the reference policy in the controlled sentiment generation environment. This is feasible as the true reward function (a sentiment classifier) is fully accessible. Conversely, outside controlled scenarios where the ground truth reward is unavailable, we compare algorithms based on their \textit{win rate} relative to a baseline policy. For this, GPT-4 serves as a stand-in for human judgment, measuring summary quality and response helpfulness in summarization and single-turn dialogue, respectively. In the summarization context, test set reference summaries define the baseline; in dialogue, it is the annotator-preferred response. While recent research indicates that language models can surpass conventional metrics as automatic evaluators \citep{Chen2023ExploringTU}, we further validate our reliance on GPT-4 with a human study discussed in Sec.~\ref{sec:human-judgments}. Results reveal strong alignment between GPT-4's assessments and those of human raters, with agreement levels between humans and GPT-4 often matching or exceeding the concordance among human annotators.

\textbf{Methods.} Besides DPO, we include several existing methods for aligning language model outputs with human preferences. For the summarization task, we test zero-shot prompting using \textbf{GPT-J} \citep{gpt-j}, while in dialogue we apply two-shot prompting with \textbf{Pythia-2.8B} \citep{biderman2023pythia}. Additionally, we assess the \textbf{SFT} approach, as well as \textbf{Preferred-FT}, where supervised fine-tuning is carried out on the selected completion $y_w$—derived from either SFT in controlled sentiment and summarization settings, or from a general LM for dialogue. Another related approach is \textbf{Unlikelihood}~\citep{welleck2019neural}, which encourages higher probability for $y_w$ while actively suppressing $y_l$, with an optional weighting parameter $\alpha\in[0,1]$ applied to the unlikelihood penalty. We further evaluate \textbf{PPO} \citep{schulman2017proximal}, employing a reward model trained from preference data, and \textbf{PPO-GT}, which acts as an oracle using the actual reward function in the controlled sentiment case. For sentiment-based experiments, we explore two PPO-GT variants: a standard off-the-shelf version \cite{leandro_von_werra_2023_7790115}, and a customized variant featuring reward normalization and enhanced hyperparameter optimization to boost performance—these modifications are also incorporated into standard PPO runs utilizing learned rewards. Lastly, the \textbf{Best of $N$} baseline is considered, whereby $N$ outputs from the SFT (or Preferred-FT in dialogue) are generated and the highest-reward response, as judged by a model trained on preference data, is selected. This method often achieves strong results but is computationally intensive at inference time, since generating $N$ samples per input can be prohibitive even for moderate $N$.

\subsection{How well can DPO optimize the RLHF objective?}

In standard RLHF frameworks, the reward maximization objective is subject to a KL constraint, which encourages maximizing reward while also limiting the extent to which the learned policy diverges from a given reference policy. Consequently, evaluations across algorithms must jointly consider both the magnitude of the reward obtained and the KL divergence; achieving a marginally greater reward at the cost of a significantly higher KL is not always optimal. In Figure~\ref{fig:frontier-tldr-main}, we present the trade-off between reward and KL across various methods within the sentiment task. For each method, we perform several independent training sessions, each with a distinct value for the policy conservativeness parameter: target KL values of $\{3,6,9,12\}$ for PPO, $\beta$ values of $\{0.05,0.1,1,5\}$, and $\alpha$ values of $\{0.05,0.1,0.5,1\}$ for the unlikelihood approach; for preferred-FT, we vary the random seed. In total, this experimental grid spans 22 distinct runs. At intervals of every 100 optimization steps up to convergence, we assess each resulting policy on a set of held-out prompts, recording both the mean true reward and the mean sequence-level KL\footnote{Specifically, this refers to the total KL divergence accumulated across all time steps.} with respect to the reference policy $\text{KL}\left(\pi\mid \mid \pi_\text{ref}\right)$. The empirical findings indicate that DPO yields a clearly superior reward-KL trade-off frontier: DPO secures the highest rewards while maintaining relatively low KL values. This performance is striking for several reasons. Notably, DPO and PPO address the same underlying optimization objective, yet DPO outperforms PPO in terms of efficiency—the reward/KL balance for DPO strictly surpasses that of PPO. Furthermore, DPO outperforms PPO even when the latter is provided access to the true reward function (PPO-GT).

\subsection{Can DPO scale to real preference datasets?}
\label{sec:dpo-real-datasets}
We proceed to assess DPO’s effectiveness in fine-tuning tasks involving summarization and single-turn dialogue. In the case of summarization, it has been shown that commonly used automatic metrics, like ROUGE, often fail to align closely with human preferences~\citep{stiennon2022learning}. Additionally, previous studies have demonstrated that language models fine-tuned with PPO on human feedback generate summaries that better satisfy human raters. Our approach compares several methods by generating summaries from the test partition of the TL;DR summarization dataset, then measuring the mean win rate against reference summaries within this test split. To evaluate all models, outputs are sampled across temperature values ranging from 0.0 to 1.0, with the resulting win rates displayed in Figure~\ref{fig:frontier-tldr-main} (right). DPO, PPO, and Preferred-FT are all fine-tuned starting from the identical GPT-J SFT base model\footnote{\url{https://huggingface.co/CarperAI/openai_summarize_tldr_sft}}. Our analysis indicates that DPO attains a win rate near 61\% at a temperature of 0.0, surpassing PPO’s approximate 57\% win rate at its own optimal temperature setting. Furthermore, DPO achieves the highest maximum win rate among all evaluated baselines, including the best-of-$N$ approach. It should be highlighted that minimal hyperparameter tuning was performed for DPO’s $\beta$, so current results might not fully reflect the approach’s upper-bound performance. DPO’s advantage is also evident in its consistent robustness to changes in sampling temperature, whereas PPO’s efficacy declines to baseline GPT-J levels at elevated temperatures. Meanwhile, Preferred-FT does not produce substantial improvements relative to the SFT baseline. Lastly, a direct human evaluation (Section~\ref{sec:human-judgments}) between DPO and PPO shows that, when comparing samples at temperature 0.25, raters favored DPO’s outputs 58\% of the time over PPO’s outputs at the same temperature.

For single-turn dialogue, we assess the various approaches using the portion of the test split from the Anthropic HH dataset \citep{bai2022training} that includes only one exchange between human and assistant. When evaluating with GPT-4, we treat the preferred completions from the test set as the gold standard to determine win rates across methods. In the absence of a standard SFT baseline for this particular task, our workflow involves beginning with a pre-trained Pythia-2.8B model, applying Preferred-FT to construct a reference model based on selected completions to maintain distributional consistency, and then conducting DPO training. Additionally, we benchmark against both the top completion from 128 Preferred-FT samples (having observed that the Best of $N$ approach saturates at 128 completions in this scenario; refer to Appendix Figure~\ref{fig:best-of-n}) and a version of the Pythia-2.8B base model prompted in a 2-shot manner. Our results show that, at optimal temperatures for each method, DPO matches or outperforms these baselines. We further analyze an RLHF system trained with PPO on the Anthropic HH dataset\footnote{\url{https://huggingface.co/reciprocate/ppo_hh_pythia-6B}} based on a reputable implementation\footnote{\url{https://github.com/CarperAI/trlx/tree/main/examples/hh}}, but could not identify a prompt or temperature configuration that would enable it to surpass the raw Pythia-2.8B performance. Drawing on our TL;DR findings and acknowledging that both Best of 128 and PPO are geared towards optimizing the same reward metric, we treat Best of 128 as an approximate stand-in for PPO-level results. Ultimately, DPO distinguishes itself as the sole method that, while remaining computationally efficient, meaningfully surpasses the preferred completion baseline within the Anthropic HH dataset, and attains equivalent or superior performance compared to the computationally intensive Best of 128 approach. Moreover, as depicted in Figure~\ref{fig:dialogue-main}, DPO rapidly reaches its optimal level of effectiveness.

\subsection{Generalization to a new input distribution}

To better assess the robustness of PPO and DPO when exposed to distributional changes, we test the policies learned from our Reddit TL;DR summarization setup on an alternative domain: news articles drawn from the CNN/DailyMail test set \citep{nallapati-etal-2016-abstractive}, utilizing the optimal sampling temperatures established for TL;DR (0 and 0.25). The findings are summarized in Table~\ref{tab:ood}. We measured the win rate of GPT-4 compared to reference summaries in the CNN/DailyMail corpus, applying the same GPT-4 (C) evaluation prompt from the Reddit TL;DR experiments, except substituting “forum post” with “news article.” In this out-of-distribution setting, DPO maintains a clear advantage over PPO. These observations offer preliminary support for the ability of DPO policies to generalize as effectively as those derived by PPO, despite DPO lacking access to the extra unlabeled Reddit TL;DR prompts leveraged by PPO.

\subsection{Validating GPT-4 judgments with human judgments}
\label{sec:human-judgments}

To test the dependability of GPT-4-based evaluation, we conducted a user study referencing TL;DR summarization experiment outcomes and employed two alternative prompts with GPT-4. The \textbf{GPT-4 (S)} (simple) prompt requests a judgment on which summary better captures key information from the post, whereas the \textbf{GPT-4 (C)} (concise) prompt incorporates an explicit preference for conciseness. We chose to examine the latter since our observations indicated that, with the simple prompt, GPT-4 favored longer, sometimes repetitive summaries, unlike human raters. Appendix~\ref{app:prompts} provides full prompt texts. Our evaluation spans three models—a top-performing (DPO, temp. 0.25), a bottom-performing (PPO, temp. 1.0), and an intermediate (SFT, temp. 0.25) method—to ensure a range of output quality. Each method is benchmarked against greedily-sampled PPO (the most successful temperature setting). Results indicate that GPT-4’s assessments align with human ratings to a similar degree as human raters concur with each other, suggesting GPT-4 is a suitable surrogate for human judgment (note that, due to limited rater availability, we only obtained multiple human ratings for DPO and PPO-1 model outputs). Notably, the \textbf{GPT-4 (C)} prompt tends to yield win rates that mirror human preferences more closely; we thus employ this prompt for headline results in Section~\ref{sec:dpo-real-datasets}. For additional specifics about the rater study—including details about the web interface and participating human volunteers—refer to Appendix~\ref{app:human-study}.

\section{Discussion}
Preference-based learning offers an effective and scalable approach to aligning and training high-performing language models. In this work, we have presented DPO, a straightforward framework that enables language models to learn from preferences without the complexities of reinforcement learning. Instead of adapting the preference learning objective to conform to a traditional RL paradigm and thereby relying on standard RL methodologies, DPO formulates a correspondence between policies in language modeling and their underlying reward functions. This enables models to align with human preferences \textit{directly} via a basic cross-entropy objective, omitting the need for reinforcement learning without sacrificing generality. Remarkably, DPO achieves competitive or superior results compared to prevalent RLHF methods such as PPO, and does so with minimal hyperparameter tuning. This substantially lowers the barrier for applying preference-based training to new language models.

\textbf{Limitations \& Future Work.} Several important avenues for further exploration emerge from our findings. One question concerns the out-of-distribution generalization of policies trained via DPO versus those optimized against explicit reward functions. Preliminary evidence indicates that DPO-trained policies generalize comparably to those trained with PPO, though a thorough analysis remains necessary. For example, it is not yet clear if leveraging self-generated labels from DPO-trained policies can harness unlabeled prompts effectively. Another research direction concerns the manifestation of reward over-optimization in the DPO paradigm and whether the minor performance reduction noted in Figure~\ref{fig:dialogue-main}-right is indicative of such effects. Our experiments, limited to models up to 6B parameters, suggest the potential for scaling DPO to much larger, state-of-the-art systems—a promising topic for future study. In terms of evaluation methodologies, we observe that GPT-4-derived win rates are sensitive to prompt phrasing; identifying optimal strategies for eliciting robust judgments from automated evaluators is an open question. Lastly, DPO's applicability extends beyond human preference alignment in language models, opening opportunities for its adoption in training generative models across diverse modalities.